\subsection{The mate lemma}
\label{subsec:mates-lemma}

In \cref{sec:2dloc} the following lemma will replace the contractibility of the walking equivalence.
A $1$-morphism from $u\colon x \to y$ to $u'\colon x' \to y'$ in $\func(\Delta^{1},\E)$ is a commutative square $\Delta^{1} \times \Delta^{1} \to \E$ with $u$, $u'$ on opposite sides.

\reviewcomment{I feel like the best way to do the mate lemma is to quote 
Haugseng's description of lax morphisms. His category of adjunctions + (co)lax
morphisms contains our category of adjunctions + strong morphisms as a
locally full subcategory, I guess? If this is true we can use his result directly.}
\begin{lem}[mates]
\label{lem:mates}
Let $\E$ be an $(\infty,2)$-category.
\begin{enumerate}
\item $\ell^{*}\colon \func(\Adj,\E) \to \func(\Delta^{1},\E)$ is a locally full inclusion. It exhibits $\func(\Adj,\E)$ as the locally full sub-$(\infty,2)$-category whose objects are the left adjoints $u\colon x \to y$, and whose $1$-morphisms from $u$ to $u'$ are the commutative squares for which the lax square
\[
\begin{tikzcd}
x \ar[r, "f_{0}"] \ar[d, "u"'] & x' \ar[d, "u'"] \ar[dl, Rightarrow, shorten=3mm, "\varphi"'] \\
y \ar[r, "f_{1}"'] & y'
\end{tikzcd}
\qquad \varphi\colon u'f_{0} \simeq f_{1}u \ ,
\]
is right Beck--Chevalley: $\rmate{\varphi}\colon f_{0}u^{R} \Rightarrow u'^{R}f_{1}$ is invertible.
\item Dually, $\rho^{*}\colon \func(\Adj,\E) \to \func(\Delta^{1},\E)$ is a locally full inclusion. It exhibits $\func(\Adj,\E)$ as the locally full sub-$(\infty,2)$-category whose objects are the right adjoints $u\colon x \to y$, and whose $1$-morphisms from $u$ to $u'$ are the commutative squares for which the lax square
\[
\begin{tikzcd}
x \ar[r, "u"] \ar[d, "f_{0}"'] & y \ar[d, "f_{1}"] \ar[dl, Rightarrow, shorten=3mm, "\psi"'] \\
x' \ar[r, "u'"'] & y'
\end{tikzcd}
\qquad \psi\colon f_{1}u \simeq u'f_{0} \ ,
\]
is left Beck--Chevalley: $\lmate{\psi}\colon u'^{L}f_{1} \Rightarrow f_{0}u^{L}$ is invertible.
\end{enumerate}
\end{lem}

\begin{proof}
We prove (1); (2) follows from the dual halves of \cref{thm:walking-adjunction,thm:adjoints-in-functor-cats}.
For $K \in \twoCat$, the map $\mapp(K,\func(\Adj,\E)) \to \mapp(K,\func(\Delta^{1},\E))$ is adjoint to
\[
\ell^{*}\colon \mapp\bigl(\Adj,\func(K,\E)\bigr) \longrightarrow \mapp\bigl(\Delta^{1},\func(K,\E)\bigr) \ ,
\]
a monomorphism with image the left adjoints of $\func(K,\E)$ by \cref{thm:walking-adjunction}. So $\ell^{*}$ is a monomorphism, and its image in $\mapp(K,\func(\Delta^{1},\E))$ consists of the functors $K \to \func(\Delta^{1},\E)$ whose adjoint $\Delta^{1} \to \func(K,\E)$ is a left adjoint. We compute it for $K = \Delta^{0}$, $\Delta^{1}$, $\wtc$.

For $K = \Delta^{0}$ it consists of the left adjoints of $\E$.
For $K = \Delta^{1}$, read a square $\Delta^{1} \times \Delta^{1} \to \E$ as a $1$-morphism $f_{0} \Rightarrow f_{1}$ of $\func(\Delta^{1},\E)$ with components $u$, $u'$; its naturality square is the lax square in (1). By \cref{thm:adjoints-in-functor-cats} it is a left adjoint if and only if $u$, $u'$ are left adjoints and this square is right Beck--Chevalley.
For $K = \wtc$, read a $2$-morphism of $\func(\Delta^{1},\E)$ between $1$-morphisms $\sigma$, $\tau$ from $u$ to $u'$ as a $1$-morphism of $\func(\wtc,\E)$ with components $u$, $u'$. Its naturality squares at the two non-identity $1$-morphisms of $\wtc$ are the lax squares of $\sigma$ and $\tau$, as restriction along $\bwtc \subset \wtc$ is compatible with the adjunction $\mapp(\wtc,\func(\Delta^{1},\E)) \simeq \mapp(\Delta^{1},\func(\wtc,\E))$. By \cref{thm:adjoints-in-functor-cats} it is a left adjoint whenever $u$, $u'$ are left adjoints and $\sigma$, $\tau$ are right Beck--Chevalley.

So every $2$-morphism of $\func(\Delta^{1},\E)$ between $1$-morphisms in the image of $\ell^{*}$ lies in the image of $\ell^{*}$. As $\ell^{*}$ is a monomorphism, this says that the square \eqref{eq:lff-square} for $G = \ell^{*}$ is a pullback (\cref{subsec:prelim-conventions}), i.e.\ that $\ell^{*}$ is locally fully faithful \cite[Prop.~2.6.5]{abellangagnahaugseng2024straightening}. Hence $\ell^{*}$ is a locally full inclusion, and the cases $K = \Delta^{0}$, $\Delta^{1}$ describe its objects and $1$-morphisms.
\end{proof}

\begin{rmk}
\label{rmk:mates-2cat}
For $2$-categories, \cref{lem:mates} is the classical fact that a morphism of adjunctions is determined by its left adjoint part, the right adjoint part being the mate \cite{kellystreet1974review}. The proof above uses only \cref{thm:walking-adjunction,thm:adjoints-in-functor-cats}.
For $(\infty,2)$-categories, the lax analogue is \cite[Cor.~4.8]{haugseng2021laxmonads}: $\ell^{*}$, resp.\ $\rho^{*}$, is an equivalence from adjunctions in $\E$ with colax, resp.\ lax, morphisms onto the full sub-$(\infty,2)$-category of $\func(\Delta^{1},\E)$ on the left, resp.\ right, adjoints. Its proof relies on the lax analogue of \cref{thm:adjoints-in-functor-cats}, which is proved in \cite[Cor.~5.2.10]{abellangagnahaugseng2024straightening} (see \cite[Rem.~5.2.11]{abellangagnahaugseng2024straightening}). Restricting to the wide locally full sub-$(\infty,2)$-category of strong transformations \cite[Cor.~2.8.12]{abellangagnahaugseng2024straightening} gives another route to \cref{lem:mates}, which requires identifying the strong morphisms of adjunctions with the Beck--Chevalley squares. We know of no reference stating \cref{lem:mates} itself.
\end{rmk}

\begin{defi}
\label{def:reflection}
An adjunction $l \dashv r$ in $\A$ is a \emph{reflection} if its counit $\varepsilon$ is invertible, and a \emph{coreflection} if its unit $\eta$ is invertible.
The left adjoint in a reflection is called a \emph{localization}, and the right adjoint in a coreflection a \emph{colocalization}.
A $1$-morphism is \emph{reflective} if it is the right adjoint of a reflection, and \emph{coreflective} if it is the left adjoint of a coreflection.
\end{defi}

By \cref{thm:walking-adjunction}, being reflective is a property of a functor $ r $ that is independant of the choice of a left adjoint,
and similarly for being coreflective.
As $r_{\natural}$ and $l_{\natural}$ on $\Hom_{\A}(t,-)$ form an adjunction with counit $\varepsilon_{\natural}$ and unit $\eta_{\natural}$, a right adjoint $r$ is reflective if and only if every $r_{\natural}$ is fully faithful, and a left adjoint $l$ is coreflective if and only if every $l_{\natural}$ is fully faithful.

Let $c\colon \wtc \to \Delta^{1}$ collapse the $2$-morphism, and let $\varepsilon,\eta\colon \wtc \to \Adj$ classify the counit $lr \Rightarrow \mathrm{id}_{+}$ and the unit $\mathrm{id}_{-} \Rightarrow rl$. The \emph{walking reflection} and \emph{walking coreflection} are the pushouts in $\twoCat$
\[
\Refl := \Adj \sqcup_{\wtc,\varepsilon} \Delta^{1} \ , \qquad \Corefl := \Adj \sqcup_{\wtc,\eta} \Delta^{1}
\]
along $c$. We write $\ell$, $\rho$ also for the composites $\Delta^{1} \to \Adj \to \Refl$, $\Corefl$.

\begin{lem}[mates for reflections]
\label{lem:mates-refl}
Let $\E$ be an $(\infty,2)$-category.
\begin{enumerate}
\item $\func(\Refl,\E) \to \func(\Adj,\E)$ is fully faithful and exhibits $\func(\Refl,\E)$ as the full sub-$(\infty,2)$-category on the reflections; likewise $\func(\Corefl,\E)$ is the full sub-$(\infty,2)$-category on the coreflections.
\item Hence for $X \in \{\Refl,\Corefl\}$, $\ell^{*}$ and $\rho^{*}\colon \func(X,\E) \to \func(\Delta^{1},\E)$ are locally full inclusions with the $1$-morphisms described in \cref{lem:mates}. The objects in the image of $\ell^{*}$, resp.\ $\rho^{*}$, are the left, resp.\ right, adjoints of reflections for $X = \Refl$, and of coreflections for $X = \Corefl$.
\end{enumerate}
\end{lem}

\begin{proof}
(1) $\mapp(\wtc,\B)$ is the space of objects $x,y \in \B$ with a morphism of $\Hom_{\B}(x,y)$, and $c^{*}$ sends $f$ to $\mathrm{id}_{f}$. As $C^{\simeq} \to \func(\Delta^{1},C)^{\simeq}$ is a monomorphism onto the equivalences for every $\infty$-category $C$, $c^{*}$ is a monomorphism onto the invertible $2$-morphisms.
Applying this to $\B = \func(K,\E)$, the map
\[
\mapp\bigl(K,\func(\Refl,\E)\bigr) \simeq \mapp\bigl(\Refl,\func(K,\E)\bigr) \longrightarrow \mapp\bigl(\Adj,\func(K,\E)\bigr) \simeq \mapp\bigl(K,\func(\Adj,\E)\bigr)
\]
is a monomorphism onto the adjunctions in $\func(K,\E)$ with invertible counit. A $2$-morphism of $\func(K,\E)$ is invertible if and only if its components at the objects of $K$ are.\todo{Reference.}
So for $K = \Delta^{0}$, $\Delta^{1}$, $\wtc$ the image consists of all $K \to \func(\Adj,\E)$ whose objects are reflections. The case of $\Corefl$ is the same.
(2) follows from (1) and \cref{lem:mates}.
\end{proof}
